% !TeX root = ../Book3.tex
\chapter*{后记}
\addcontentsline{toc}{chapter}{后记}
从一个切空间中的内积开始，本卷逐渐走到距离空间、弱梯度和奇异极限。贯穿这一过程的，是怎样把局部计算变成整体控制：沿一条测地线的二阶变分可以约束直径，曲率积分可以给出拓扑数，能量估计可以保证一列解存在可用的极限。每一次转换都需要找到能够保留的信息，并选择适合承载它的对象。

\ProseParagraphBreak
同一个度量在不同问题中承担不同职责。它规定长度和体积，也确定函数空间中的范数、方程的主符号及能量的尺度。研究一个退化序列时，距离、体积测度与微分算子未必以同样的方式收敛；研究一个变分问题时，弱收敛、能量下半连续和正则性也分别解决不同的障碍。重读时可以选一个已经算清的模型，逐项追踪这些结构如何变化，而不急于把所有收敛归入同一个名称。

\ProseParagraphBreak
圆球、平坦环面、双曲商和曲面映射反复出现，为这样的追踪提供了具体起点。圆球把曲率、特征值、临界能量和自相似流联系起来；平坦环面显示局部曲率不能独自决定整体形状；双曲尖点区分了有限体积与紧致性；曲面上的气泡则说明有界能量仍然允许尺度集中。这些计算使一般定理中的条件具有可以辨认的几何含义。

\ProseParagraphBreak
本卷为椭圆估计、紧性、变分和局部模型保留了较多证明细节。一个截断函数为何这样选择，一项误差怎样吸收，一个规范条件怎样恢复椭圆性，常常决定方法能否用于另一个问题。继续学习时，可以尝试改变一个条件，检查证明最先在哪一步失效，再寻找替代估计。这样的工作也能帮助区分结论本身的必要条件与某种证明方法所需的便利条件。

\ProseParagraphBreak
后续阅读已有几条具体方向。关心曲率与拓扑，可以继续研究特殊完整群、分裂与刚性；关心方程，可以深入临界紧性、几何流奇点及规范理论；关心面积问题，可以比较整数流与变分流形的存在性、边界和正则性。进入一般结构定理以前，先掌握本卷已经展开的模型及其归一化，再对照所引文献的对象、维数和收敛方式，会更容易看清新工具增加了什么。

\ProseParagraphBreak
第四卷将把联络、椭圆分析和变分方法放到复结构与辛结构中，研究相容度量、全纯对象和 Hamilton 系统；第五卷则用代数几何的语言继续讨论局部方程、上同调、相交与模问题。它们与本卷共享一些几何对象，也带来各自的定义和分类方法。能够从一个明确的问题出发，选对对象与估计，并把局部结论接到整体几何，便有了继续探索这些联系的基础。
